\chapter{Υλοποίηση της αποκεντρωμένης εφαρμογής}
\label{ch:dapp}

Το κεφάλαιο περιγράφει την εφαρμογή ιστού μέσω της οποίας οι χρήστες αλληλεπιδρούν με το
συμβόλαιο. Οι πρώτες ενότητες παρουσιάζουν την τελική υλοποίηση· η
Ενότητα~\ref{sec:difficulties} καταγράφει τις δυσκολίες που προέκυψαν και τις προσεγγίσεις
που δοκιμάσθηκαν, στο πνεύμα του αντίστοιχου κεφαλαίου της \cite{hadjichristofi2022}· και
η Ενότητα~\ref{sec:verification} περιγράφει πώς επαληθεύθηκε ότι το σύνολο λειτουργεί.

\section{Επισκόπηση και τεχνολογίες}
\label{sec:stack}

Η εφαρμογή είναι μια \en{single-page application} που εκτελείται εξ ολοκλήρου στον
φυλλομετρητή. Δεν διαθέτει δικό της διακομιστή: τα αρχεία της είναι στατικά, και κάθε
δεδομένο που εμφανίζει προέρχεται απευθείας από έναν πλήρη κόμβο του \en{Sui}. Αποτελείται
από περίπου 4.200 γραμμές \en{TypeScript} και \en{CSS} σε 57 αρχεία. Ο Πίνακας~\ref{tab:stack}
παραθέτει τις τεχνολογίες που χρησιμοποιήθηκαν.

\begin{table}[H]
\centering
\caption{Τεχνολογίες της εφαρμογής}
\label{tab:stack}
\small
\begin{tabular}{llp{7.6cm}}
\toprule
Τεχνολογία & Έκδοση & Ρόλος \\
\midrule
\en{TypeScript} \cite{typescript} & 5.9 & Γλώσσα· στατικός έλεγχος τύπων σε όλο τον κώδικα. \\
\en{React} \cite{react} & 19 & Βιβλιοθήκη διεπαφής χρήστη. \\
\en{Vite} \cite{vite} & 6 & Διακομιστής ανάπτυξης και εργαλείο παραγωγής. \\
\en{Sui TypeScript SDK} \cite{sui-ts-sdk} & 2.31 & Κατασκευή συναλλαγών, πελάτες \en{gRPC} και \en{JSON-RPC}. \\
\en{dApp Kit React} \cite{sui-dapp-kit} & 2.1 & Σύνδεση πορτοφολιού, υπογραφή, \en{React hooks}. \\
\en{TanStack Query} \cite{tanstack-query} & 5 & Ανάκτηση, προσωρινή αποθήκευση και ακύρωση δεδομένων. \\
\en{Tailwind CSS} \cite{tailwind} & 4 & Μορφοποίηση με βάση \en{design tokens}. \\
\en{Radix Primitives} \cite{radix} & 1.1 & Προσβάσιμα παράθυρα διαλόγου, καρτέλες, πλαίσια επιλογής. \\
\en{React Router} & 7 & Δρομολόγηση και κατάσταση φίλτρων στη διεύθυνση \en{URL}. \\
\en{Geist}, \en{Phosphor Icons} & 5.3, 2.1 & Γραμματοσειρές (τοπικά φιλοξενούμενες) και εικονίδια. \\
\en{sonner} & 1.7 & Ειδοποιήσεις (\en{toasts}) με σύνδεσμο στον εξερευνητή. \\
\bottomrule
\end{tabular}
\end{table}

Ο κώδικας οργανώνεται σε στρώματα που αντιστοιχούν στο Σχήμα~\ref{fig:architecture}:

\begin{lstlisting}[numbers=none,caption={Δομή του καταλόγου \code{app/src}},label={lst:app-tree}]
src/
  dapp-kit.ts          # στιγμιότυπο dApp Kit, δίκτυα
  config/              # διεύθυνση πακέτου, modules, τύποι, κατάλογος NFT
  lib/                 # πελάτης Sui, ανάλυση αντικειμένων, εξερευνητές ανά δίκτυο
  hooks/               # ένα hook ανά ερώτημα και ανά συναλλαγή, faucet, υπόλοιπο
  components/
    ui/                # γενικά στοιχεία: Button, Dialog, Tabs, ...
    nft/               # κάρτα NFT, εικόνα, σπανιότητα, φίλτρο σπανιότητας
    bid/               # υποβολή προσφοράς, εισερχόμενες, γραμμή προσφοράς
    layout/            # γραμμή πλοήγησης, υπόλοιπο και faucet
    TxToast.tsx        # ειδοποίηση συναλλαγής με σύνδεσμο στον εξερευνητή
    ExplorerLink.tsx   # σύνδεσμος αντικειμένου ή συναλλαγής στον εξερευνητή
  pages/               # οι πέντε σελίδες
  types/               # τύποι NFT και προσφοράς
\end{lstlisting}

Ο κανόνας που τηρείται είναι ότι οι σελίδες και τα στοιχεία διεπαφής δεν γνωρίζουν τίποτε
για το \en{Sui}. Καλούν \en{hooks}, και οι \en{hooks} είναι το μόνο στρώμα που
κατασκευάζει συναλλαγές ή διαβάζει την αλυσίδα.

\section{Σύνδεση με το δίκτυο και το πορτοφόλι}
\label{sec:connection}

\subsection{Το στιγμιότυπο dApp Kit}

Η σύνδεση με πορτοφόλια και δίκτυα διαμορφώνεται σε ένα σημείο:

\lstinputlisting[language=TypeScript,firstline=4,lastline=9,firstnumber=4,
  caption={[Το στιγμιότυπο \en{dApp Kit}]Το στιγμιότυπο \en{dApp Kit} (\code{dapp-kit.ts})},label={lst:dappkit}]{\appsrc/dapp-kit.ts}

Το \en{dApp Kit} ανακαλύπτει αυτόματα κάθε πορτοφόλι που υλοποιεί το \en{Wallet Standard},
αποκαθιστά την προηγούμενη σύνδεση όταν ο χρήστης επιστρέφει, και εκθέτει \en{hooks} για
τον τρέχοντα λογαριασμό και τον τρέχοντα πελάτη. Η επιλογή \code{enableBurnerWallet}
ενεργοποιεί, μόνο σε περιβάλλον ανάπτυξης, ένα προσωρινό πορτοφόλι μέσα στον φυλλομετρητή
που υπογράφει χωρίς επιβεβαίωση. Η δυνατότητα αυτή αποδείχθηκε καθοριστική για τον
αυτοματοποιημένο έλεγχο της διεπαφής (Ενότητα~\ref{sec:verification}).

\subsection{Πελάτης ανά δίκτυο}

Η εφαρμογή δηλώνει τρία δίκτυα, \en{localnet}, \en{devnet} και \en{testnet}
(Ενότητα~\ref{sec:sui-networks}), και ένα από αυτά ως προεπιλογή· η έκδοση ανάπτυξης
προεπιλέγει το τοπικό δίκτυο και η τελική έκδοση το \en{testnet}
(Ενότητα~\ref{sec:dev-flow}). Το \en{SDK} προσφέρει περισσότερους από έναν πελάτες,
ανάλογα με το πρωτόκολλο επικοινωνίας με τον κόμβο. Ο συνιστώμενος είναι ο \code{SuiGrpcClient}
\cite{sui-grpc}· ο παλαιότερος \code{SuiJsonRpcClient} διατηρείται για συμβατότητα. Και οι
δύο υλοποιούν μια κοινή διεπαφή, το \en{Core API}, προσβάσιμη ως \code{client.core}. Η
εφαρμογή επιλέγει πελάτη ανά δίκτυο:

\lstinputlisting[language=TypeScript,firstline=5,lastline=23,firstnumber=5,
  caption={[Επιλογή πελάτη ανά δίκτυο]Επιλογή πελάτη ανά δίκτυο (\code{lib/suiClient.ts})},label={lst:client}]{\appsrc/lib/suiClient.ts}

Ο λόγος για τον οποίο το τοπικό δίκτυο χρησιμοποιεί διαφορετικό πελάτη εξηγείται στην
Ενότητα~\ref{sec:listevents}. Το σημαντικό εδώ είναι ότι ο υπόλοιπος κώδικας δεν το
γνωρίζει: όλες οι αναγνώσεις γίνονται μέσω του \code{client.core}, και η αλλαγή πελάτη για
ένα δίκτυο είναι αλλαγή μίας γραμμής.

\section{Ανάγνωση της κατάστασης χωρίς indexer}
\label{sec:reading}

Η εφαρμογή χρειάζεται τρία είδη πληροφορίας: ποια \en{NFTs} κατέχει ο χρήστης, ποια
\en{NFTs} υπάρχουν συνολικά, και ποιες προσφορές είναι ενεργές. Μόνο το πρώτο απαντάται
με ένα ερώτημα σε πλήρη κόμβο.

\subsection{NFTs του χρήστη}

Ο κόμβος ευρετηριάζει τα αντικείμενα ανά ιδιοκτήτη. Η μέθοδος
\code{listOwnedObjects}, με φίλτρο τον τύπο \code{MarketplaceNFT}, επιστρέφει σελιδοποιημένα
τα \en{NFTs} μιας διεύθυνσης. Τα \en{NFTs} που ο χρήστης έχει δεσμεύσει σε προσφορές δεν
περιλαμβάνονται, και ορθώς: ανήκουν στο αντικείμενο \code{Bid}. Η σελίδα «\en{My NFTs}»
τα εμφανίζει παρ' όλα αυτά, με την ένδειξη «\en{In a bid}», αντλώντας τα από τις
δεσμεύσεις των ενεργών προσφορών του χρήστη (ερώτημα \code{useMyBids},
Πίνακας~\ref{tab:queries}), ώστε ο χρήστης να βλέπει ολόκληρη την περιουσία του και να μπορεί να
ακυρώσει μια προσφορά από το ίδιο το \en{NFT}.

\subsection{Ανακάλυψη μέσω συμβάντων}

Για τα υπόλοιπα δύο, η εφαρμογή ξεκινά από τα συμβάντα
(Ενότητα~\ref{sec:events}). Η βοηθητική \code{listAllEvents} διαβάζει όλα τα συμβάντα ενός
τύπου, από το νεότερο προς το παλαιότερο:

\lstinputlisting[language=TypeScript,firstline=14,lastline=34,firstnumber=14,
  caption={[Ανάγνωση όλων των συμβάντων ενός τύπου]Ανάγνωση όλων των συμβάντων ενός τύπου (\code{lib/bidEnriched.ts})},label={lst:events}]{\appsrc/lib/bidEnriched.ts}

Η σελιδοποίηση δεν είναι τυπικότητα. Στη διεπαφή \en{gRPC}, ο κόμβος περιορίζει το τμήμα
της αλυσίδας που σαρώνει ανά αίτημα, και μια σελίδα μπορεί να επιστραφεί κενή ενώ
υπάρχουν παλαιότερα συμβάντα. Ο βρόχος συνεχίζει όσο ο κόμβος δηλώνει ότι υπάρχει επόμενη
σελίδα.

Το Σχήμα~\ref{fig:pipeline} δείχνει πώς από τα συμβάντα προκύπτει η λίστα των
εισερχόμενων προσφορών ενός χρήστη, και το Απόσπασμα~\ref{lst:received} την υλοποίηση.

\begin{figure}[ht]
\centering
\begin{tikzpicture}[
  font=\footnotesize, >=Stealth,
  stp/.style={draw, rounded corners=3pt, minimum width=2.55cm, minimum height=1.25cm, align=center, fill=gray!8, text width=2.35cm}]
  \node[stp] (a) at (0,0) {συμβάντα\\\code{BidPlaced}};
  \node[stp, right=0.45cm of a] (b) {φίλτρο:\\όχι δικές μου,\\με στόχο};
  \node[stp, right=0.45cm of b] (c) {αντικείμενα\\\code{Bid}\\(όσα υπάρχουν)};
  \node[stp, right=0.45cm of c] (d) {\en{NFTs}-στόχοι·\\φίλτρο: κάτοχος\\= εγώ};
  \node[stp, right=0.45cm of d] (e) {δεσμευμένα\\\en{NFTs}\\(\en{DOF})};
  \foreach \x/\y in {a/b,b/c,c/d,d/e} \draw[->] (\x) -- (\y);
  \node[font=\scriptsize, below=0.1cm of a] {\code{listEvents}};
  \node[font=\scriptsize, below=0.1cm of c] {\code{getObjects}};
  \node[font=\scriptsize, below=0.1cm of d] {\code{getObjects}};
  \node[font=\scriptsize, below=0.1cm of e] {\code{listDynamicFields}};
\end{tikzpicture}
\caption{Ανακάλυψη των εισερχόμενων προσφορών ενός χρήστη από δεδομένα της αλυσίδας.}
\label{fig:pipeline}
\end{figure}

\lstinputlisting[language=TypeScript,firstline=15,lastline=71,firstnumber=15,
  caption={[Το ερώτημα εισερχόμενων προσφορών]Το ερώτημα εισερχόμενων προσφορών (\code{hooks/useReceivedBids.ts})},label={lst:received}]{\appsrc/hooks/useReceivedBids.ts}

Τέσσερα σημεία αξίζουν σχολιασμό.

\emph{Ο προσφέρων είναι ο αποστολέας του συμβάντος.} Το ωφέλιμο φορτίο του
\code{BidPlaced} δεν περιέχει πεδίο \code{bidder}· ο κόμβος καταγράφει τον αποστολέα
κάθε συμβάντος στα μεταδεδομένα του, και το \en{SDK} τον εκθέτει ως \code{e.sender}. Το
φίλτρο «όχι δικές μου» και το αντίστροφό του στο \code{useMyBids} συγκρίνουν αυτό το πεδίο.

\emph{Οι τερματισμένες προσφορές φιλτράρονται από την ύπαρξη του αντικειμένου.} Το συμβάν
\code{BidPlaced} παραμένει για πάντα στο ιστορικό. Αν η προσφορά έγινε αποδεκτή, απορρίφθηκε
ή ακυρώθηκε, το αντικείμενο \code{Bid} έχει διαγραφεί, και η \code{getObjects} επιστρέφει
σφάλμα για το συγκεκριμένο \en{ID} αντί για αντικείμενο. Η \code{parseBidObject} μετατρέπει
τα σφάλματα σε \code{null}, και έτσι απομένουν μόνο οι ενεργές προσφορές. Δεν χρειάζεται
να διαβασθούν τα συμβάντα τερματισμού.

\emph{Ο αποδέκτης προσδιορίζεται από τον τρέχοντα κάτοχο του στόχου.} Το συμβάν δεν
αναφέρει σε ποιον απευθύνεται η προσφορά, διότι η προσφορά στοχεύει αντικείμενο και όχι
πρόσωπο. Η εφαρμογή ανακτά τα στοχευόμενα \en{NFTs} και κρατά τις προσφορές των οποίων ο
στόχος ανήκει \emph{τώρα} στον χρήστη. Αν ένα \en{NFT} αλλάξει κάτοχο, οι προσφορές του
ακολουθούν αυτόματα τον νέο.

\emph{Τα δεσμευμένα \en{NFTs} διαβάζονται ως δυναμικά πεδία.} Η \code{fetchOfferedNfts}
απαριθμεί τα δυναμικά πεδία της προσφοράς και, επειδή πρόκειται για πεδία αντικειμένων,
λαμβάνει απευθείας τα \en{IDs} των παιδιών:

\lstinputlisting[language=TypeScript,firstline=121,lastline=142,firstnumber=121,
  caption={[Ανάγνωση των δεσμευμένων \en{NFTs} μιας προσφοράς]Ανάγνωση των δεσμευμένων \en{NFTs} μιας προσφοράς (\code{lib/bidEnriched.ts})},label={lst:offered}]{\appsrc/lib/bidEnriched.ts}

Με τον ίδιο τρόπο η σελίδα μιας συλλογής προκύπτει από τα συμβάντα \code{NftMinted}. Εκεί
το φίλτρο είναι ο ιδιοκτήτης: εμφανίζονται μόνο τα \en{NFTs} που ανήκουν σε διεύθυνση. Όσα
είναι δεσμευμένα σε προσφορά ανήκουν σε αντικείμενο και παραλείπονται, ώστε να μην μπορεί
κάποιος να υποβάλει προσφορά για ένα \en{NFT} που ο κάτοχός του έχει ήδη προσφέρει αλλού.
Η σελίδα προσφέρει επιπλέον φίλτρο σπανιότητας, πολλαπλής επιλογής, του οποίου η
κατάσταση ζει στη διεύθυνση \en{URL} (\code{?rarity=common\&rarity=rare}): η επιλογή
επιβιώνει σε επαναφόρτωση, μοιράζεται ως σύνδεσμος, και ένα ίδιο στοιχείο εξυπηρετεί και
τις τρεις συλλογές.

\subsection{Μεταδεδομένα και προσωρινή αποθήκευση}

Για κάθε αντικείμενο ζητούνται δύο όψεις: η απόδοση \en{Display} και η αναπαράσταση
\en{JSON} των πεδίων. Το όνομα, η εικόνα και τα λοιπά στοιχεία λαμβάνονται από το
\en{Display} όταν είναι διαθέσιμο και, διαφορετικά, από τα πεδία. Η εφεδρική αυτή διαδρομή
αποδείχθηκε απαραίτητη (Ενότητα~\ref{sec:other-difficulties}).

Όλα τα ερωτήματα εκτελούνται μέσω του \en{TanStack Query}, που τα αποθηκεύει προσωρινά με
ένα κλειδί. Ο Πίνακας~\ref{tab:queries} δείχνει τα ερωτήματα και από ποιες συναλλαγές
ακυρώνεται το καθένα.

\begin{table}[H]
\centering
\caption{Ερωτήματα της εφαρμογής και ακύρωση της προσωρινής μνήμης}
\label{tab:queries}
\small
\begin{tabular}{llp{6.4cm}}
\toprule
\en{Hook} & Κλειδί & Ακυρώνεται μετά από \\
\midrule
\code{useOwnedNFTs} & \code{ownedNFTs} & κάθε συναλλαγή \\
\code{useMarketplaceNFTs} & \code{marketplaceNFTs} & δημιουργία, υποβολή, αποδοχή, ακύρωση \\
\code{useReceivedBids} & \code{receivedBids} & αποδοχή, απόρριψη \\
\code{useMyBids} & \code{myBids} & υποβολή, αποδοχή, ακύρωση \\
\code{useSuiBalance} & \code{suiBalance} & κάθε συναλλαγή, \en{faucet}· και περιοδικά ανά 15~s \\
\bottomrule
\end{tabular}
\end{table}

\section{Σύνθεση και εκτέλεση συναλλαγών}
\label{sec:transactions}

\subsection{Κατασκευή των PTBs}

Κάθε ενέργεια του χρήστη αντιστοιχεί σε ένα \en{hook} που κατασκευάζει ένα \en{PTB}. Η
υποβολή προσφοράς είναι το πιο σύνθετο:

\lstinputlisting[language=TypeScript,firstline=20,lastline=32,firstnumber=20,
  caption={[Το \en{PTB} της υποβολής προσφοράς]Το \en{PTB} της υποβολής προσφοράς (\code{hooks/usePlaceBid.ts})},label={lst:placebid-ts}]{\appsrc/hooks/usePlaceBid.ts}

Ο κώδικας αντιστοιχεί ένα προς ένα στο Σχήμα~\ref{fig:ptb}. Η \code{tx.object(id)} δηλώνει
ένα αντικείμενο ως είσοδο με το \en{ID} του, και το \en{SDK} συμπληρώνει κατά την κατασκευή
την έκδοση και το είδος ιδιοκτησίας του. Η \code{tx.gas} αναφέρεται στο νόμισμα από το οποίο
θα πληρωθεί το \en{gas}, και το προσφερόμενο ποσό αποκόπτεται από αυτό, ώστε ο χρήστης να
μην χρειάζεται να επιλέξει νομίσματα. Τα ορίσματα τύπου δίνονται ρητά στο
\code{typeArguments}. Στην εφαρμογή επίδειξης και οι δύο τύποι είναι ο
\code{MarketplaceNFT}, αφού οι τρεις συλλογές της μοιράζονται έναν τύπο \en{Move}· για
προσφορά μεταξύ διαφορετικών τύπων αρκεί να αλλάξουν οι δύο αυτές τιμές.

Ο Πίνακας~\ref{tab:ptbs} συνοψίζει τα \en{PTBs} όλων των ροών.

\begin{table}[H]
\centering
\caption{Τα \en{PTBs} της εφαρμογής}
\label{tab:ptbs}
\small
\begin{tabular}{lp{7.6cm}c}
\toprule
Ροή & Εντολές & Υπογραφές \\
\midrule
Δημιουργία $n$ \en{NFTs} & $n \times$ \code{MoveCall mint}, \code{TransferObjects} & 1 \\
Υποβολή προσφοράς & \code{MakeMoveVec}, \code{SplitCoins}, \code{MoveCall place\_bid} & 1 \\
Αποδοχή & \code{MakeMoveVec}, \code{MoveCall accept\_bid} & 1 \\
Απόρριψη & \code{MoveCall reject\_bid} & 1 \\
Ακύρωση & \code{MoveCall cancel\_bid} & 1 \\
\bottomrule
\end{tabular}
\end{table}

Κάθε ροή απαιτεί μία υπογραφή (ΜΛΑ8). Η δημιουργία πολλών \en{NFTs} γίνεται σε μία
συναλλαγή: η \code{mint} επιστρέφει κάθε αντικείμενο ως αποτέλεσμα εντολής, και μία τελική
εντολή \code{TransferObjects} παραδίδει ολόκληρη την παρτίδα στον αποστολέα
(Απόσπασμα~\ref{lst:mint-ts}). Η ίδια τεχνική θα επέτρεπε δημιουργία και προσφορά στην
ίδια συναλλαγή.

\lstinputlisting[language=TypeScript,firstline=21,lastline=34,firstnumber=21,
  caption={[Δημιουργία παρτίδας \en{NFTs}]Δημιουργία παρτίδας \en{NFTs} σε ένα \en{PTB} (\code{hooks/useMintNFT.ts})},label={lst:mint-ts}]{\appsrc/hooks/useMintNFT.ts}

\subsection{Εκτέλεση, αναμονή και ανανέωση}

Όλα τα \en{hooks} συναλλαγών μοιράζονται έναν βοηθητικό μηχανισμό εκτέλεσης:

\lstinputlisting[language=TypeScript,firstline=7,lastline=30,firstnumber=7,
  caption={[Κοινός μηχανισμός εκτέλεσης συναλλαγών]Κοινός μηχανισμός εκτέλεσης συναλλαγών (\code{hooks/useTransactionMutation.ts})},label={lst:mutation}]{\appsrc/hooks/useTransactionMutation.ts}

Ο μηχανισμός κάνει τρία πράγματα που είναι εύκολο να παραλειφθούν.

Πρώτον, ελέγχει το \emph{είδος} του αποτελέσματος. Μια συναλλαγή μπορεί να καταγραφεί στην
αλυσίδα και να έχει αποτύχει, για παράδειγμα επειδή μια \code{assert!} του συμβολαίου δεν
ικανοποιήθηκε. Το \en{SDK} επιστρέφει τότε αποτέλεσμα τύπου \code{FailedTransaction} αντί
να εγείρει εξαίρεση, και ένας κώδικας που δεν το ελέγχει θα ανακοίνωνε επιτυχία. Στην
πράξη, οι περισσότερες τέτοιες αποτυχίες εντοπίζονται νωρίτερα, διότι το \en{SDK}
προσομοιώνει τη συναλλαγή πριν ζητήσει υπογραφή.

Δεύτερον, περιμένει να \emph{ευρετηριασθεί} η συναλλαγή πριν επιστρέψει. Η ολοκλήρωση μιας
συναλλαγής δεν συνεπάγεται ότι ο κόμβος απαντά ήδη στα ερωτήματα με τη νέα κατάσταση.
Χωρίς την αναμονή, η άμεση ανανέωση των δεδομένων θα μπορούσε να επιστρέψει την παλιά
εικόνα.

Τρίτον, επιτρέπει σε κάθε \en{hook} να ακυρώσει, μετά την επιτυχία, ακριβώς τα ερωτήματα
που επηρεάζονται (Πίνακας~\ref{tab:queries}), ώστε η διεπαφή να ενημερώνεται χωρίς
επαναφόρτωση. Το υπόλοιπο \en{SUI} ανανεώνεται μετά από κάθε συναλλαγή ανεξαρτήτως
έκβασης, αφού και μια αποτυχημένη συναλλαγή καταναλώνει \en{gas}.

Κάθε επιτυχής ενέργεια επιστρέφει το ψηφιακό αποτύπωμα (\en{digest}) της συναλλαγής, και
η ειδοποίηση που βλέπει ο χρήστης το μετατρέπει σε σύνδεσμο «\en{View on \ldots}» προς
έναν εξερευνητή αλυσίδας. Ο εξερευνητής επιλέγεται ανά δίκτυο
(\code{lib/explorer.ts}): \en{Polymedia Explorer} για το τοπικό δίκτυο, επειδή είναι ο
μόνος που μπορεί να διαβάσει έναν κόμβο στον υπολογιστή του χρήστη, \en{SuiVision} για το
\en{testnet} και \en{Suiscan} για το \en{devnet}. Ο ίδιος πίνακας τροφοδοτεί τα εικονίδια
δίπλα σε κάθε διεύθυνση και \en{NFT}.

Τα σφάλματα παρουσιάζονται στον χρήστη με κατανοητή διατύπωση. Μια βοηθητική συνάρτηση
αναγνωρίζει τις συνηθέστερες περιπτώσεις (απόρριψη στο πορτοφόλι, ανεπαρκές υπόλοιπο) και
αντιστοιχίζει τους επτά κωδικούς σφαλμάτων του συμβολαίου σε μηνύματα, ώστε ένα
«\en{abort code: 0}» να εμφανίζεται ως «\en{Only the bidder can do that}».

\section{Υλοποίηση της διεπαφής}
\label{sec:ui-impl}

Οι αρχές σχεδίασης της Ενότητας~\ref{sec:ui-design} υλοποιήθηκαν σε τρία επίπεδα, και η
διεπαφή πέρασε στη συνέχεια από έναν δεύτερο σχεδιαστικό γύρο.

\paragraph{Design tokens.} Τα χρώματα, οι γραμματοσειρές, οι σκιές, οι καμπύλες κίνησης
και οι ακτίνες καμπυλότητας των μακετών ορίζονται μία φορά, ως μεταβλητές θέματος του
\en{Tailwind}. Από αυτές παράγονται οι κλάσεις που χρησιμοποιούν όλα τα στοιχεία. Οι τρεις βαθμίδες
σπανιότητας ορίζονται ως κλάσεις με μεταλλική διαβάθμιση και χρησιμοποιούνται από τρία
στοιχεία: ετικέτα, κουκκίδα και πλαίσιο.

\paragraph{Γενικά στοιχεία.} Ένα μικρό σύνολο στοιχείων, όπως \code{Button},
\code{Dialog}, \code{Tabs} και \code{Badge}, ενσωματώνει τις επαναλαμβανόμενες αποφάσεις.
Το \code{Button} διαθέτει ιδιότητα \code{loading}, που εμφανίζει ένδειξη προόδου και
απενεργοποιεί το πάτημα· έτσι η «ορατή εκκρεμότητα» υλοποιείται με τον ίδιο τρόπο παντού.
Στη σελίδα δημιουργίας, όσο ένα \en{mint} εκκρεμεί, οι υπόλοιπες κάρτες παραμένουν
ενεργές: ένα δεύτερο πάτημα μπαίνει σε ουρά («\en{Queued}») και εκτελείται μόλις
ολοκληρωθεί το πρώτο, διότι δύο ταυτόχρονες συναλλαγές του ίδιου πορτοφολιού θα
διεκδικούσαν το ίδιο νόμισμα \en{gas}.
Το \code{Dialog} έχει τρεις παραλλαγές: κεντρικό παράθυρο, φύλλο που ανασύρεται από το
κάτω μέρος της οθόνης στα κινητά, και πλήρη οθόνη για τη σύνθεση προσφοράς.

\paragraph{Στοιχεία πεδίου.} Η κάρτα \en{NFT}, η γραμμή προσφοράς, τα «χάπια» σύνοψης
και τα πλαίσια «δίνεις / παίρνεις» είναι τα στοιχεία που συνθέτουν τις σελίδες. Η γραμμή
προσφοράς δέχεται την κατεύθυνση ως παράμετρο και αντιστρέφει τις δύο πλευρές, ώστε η ίδια
προσφορά να παρουσιάζεται σωστά και στον προσφέροντα και στον κάτοχο (ΛΑ10).

Η σύνθεση προσφοράς κάνει το βήμα του \en{SUI} ρητά προαιρετικό, με κουμπί
«\en{Skip}», και η γραμμή πλοήγησης δείχνει το υπόλοιπο του πορτοφολιού και ένα κουμπί
\en{faucet} που ζητά δοκιμαστικά \en{SUI} από το \en{faucet} του τρέχοντος δικτύου, ώστε
ένα νέο προσωρινό πορτοφόλι να μπορεί να χρησιμοποιηθεί χωρίς τη γραμμή εντολών.

\paragraph{Δεύτερος σχεδιαστικός γύρος.} Μετά την υλοποίηση των μακετών, η διεπαφή
πέρασε από έναν έλεγχο με βάση ένα δημοσιευμένο σύνολο κανόνων σχεδίασης
διεπαφών~\cite{taste-skill}. Ο έλεγχος δεν άλλαξε τη δομή των σελίδων, που είχε ήδη
εγκριθεί, αλλά τα υλικά τους: οι γραμματοσειρές \en{Inter} και \en{DM Mono}, που
φορτώνονταν από εξωτερική υπηρεσία, αντικαταστάθηκαν από τις \en{Geist Sans} και
\en{Geist Mono}, φιλοξενούμενες τοπικά· τα εικονίδια ενοποιήθηκαν σε μία οικογένεια
(\en{Phosphor}) με ένα πάχος γραμμής· οι κάρτες έχασαν το γκρίζο περίγραμμα υπέρ μιας
λεπτής γραμμής και σκιάς χρωματισμένης στην απόχρωση της μάρκας, με ομόκεντρες ακτίνες
πλαισίου και εικόνας· όλες οι μεταβάσεις απέκτησαν κοινή καμπύλη κίνησης και τα κουμπιά
απόκριση στο πάτημα· και αποκαλύφθηκε ότι τα παράθυρα διαλόγου άνοιγαν ακαριαία, επειδή
οι κλάσεις κίνησης αναφέρονταν σε ένα πρόσθετο που δεν είχε εγκατασταθεί. Όλες οι
κινήσεις σέβονται τη ρύθμιση \code{prefers-reduced-motion}.

Η διεπαφή προσαρμόζεται σε πλάτη από 360 εικονοστοιχεία. Στα κινητά η πλοήγηση μεταφέρεται
σε δεύτερη γραμμή, τα πλέγματα γίνονται δίστηλα, η σελίδα δημιουργίας εμφανίζει μία
συλλογή τη φορά, και οι ενέργειες των προσφορών εμφανίζονται μόνο στην ανεπτυγμένη γραμμή,
με κουμπιά ύψους κατάλληλου για αφή.

\section{Δυσκολίες και εναλλακτικές προσεγγίσεις}
\label{sec:difficulties}

\subsection{Το ListEvents στο τοπικό δίκτυο}
\label{sec:listevents}

Η πρώτη εκδοχή της εφαρμογής χρησιμοποιούσε τον \code{SuiGrpcClient} σε όλα τα δίκτυα και
μεταγλωττιζόταν χωρίς σφάλματα. Το πρόβλημα φάνηκε μόνο όταν εκτελέσθηκε ένα σενάριο επαλήθευσης σε
πραγματικό τοπικό δίκτυο: η δημιουργία \en{NFTs}, η ανάγνωση αντικειμένων και η εκτέλεση
συναλλαγών λειτουργούσαν μέσω \en{gRPC}, αλλά κάθε ερώτημα συμβάντων αποτύγχανε με το
σφάλμα \code{UNIMPLEMENTED}. Ο κόμβος που εκκινεί η εντολή \code{sui start}, στην έκδοση
1.67 της αλυσίδας εργαλείων, δεν υλοποιεί τη μέθοδο \code{ListEvents} της υπηρεσίας
\en{gRPC}. Εφόσον
τρία από τα ερωτήματα της εφαρμογής ξεκινούν από συμβάντα (Πίνακας~\ref{tab:queries}), η
εφαρμογή θα εμφανιζόταν άδεια, χωρίς κανένα σφάλμα μεταγλώττισης να έχει προειδοποιήσει.

Εξετάσθηκαν τρεις λύσεις:
\begin{enumerate}[nosep]
  \item Εκκίνηση του τοπικού δικτύου μαζί με \en{indexer} και υπηρεσία \en{GraphQL}.
        Απαιτεί βάση δεδομένων και αναιρεί την απλότητα του περιβάλλοντος ανάπτυξης.
  \item Χρήση του \en{JSON-RPC} για όλα τα δίκτυα. Λειτουργεί, αλλά δεσμεύει την
        εφαρμογή σε διεπαφή που έχει χαρακτηρισθεί παρωχημένη.
  \item Επιλογή πελάτη ανά δίκτυο, με όλες τις αναγνώσεις μέσω του \en{Core API}, το οποίο
        υλοποιούν και οι δύο πελάτες.
\end{enumerate}
Επιλέχθηκε η τρίτη (Απόσπασμα~\ref{lst:client}). Το τοπικό δίκτυο χρησιμοποιεί
\en{JSON-RPC}, τα δημόσια δίκτυα \en{gRPC}, και όταν ο τοπικός κόμβος αποκτήσει τη μέθοδο
που λείπει, η αλλαγή θα είναι μία γραμμή. Το δίδαγμα είναι γενικότερο: ο στατικός έλεγχος
τύπων επαληθεύει ότι ο κώδικας χρησιμοποιεί σωστά μια διεπαφή, όχι ότι ο διακομιστής στην
άλλη άκρη την υλοποιεί.

\subsection{Λοιπές δυσκολίες}
\label{sec:other-difficulties}

\paragraph{Display μέσω gRPC.} Η διεπαφή \en{gRPC} επιστρέφει μεταδεδομένα προβολής μόνο
για τη νεότερη μορφή του προτύπου \en{Display}. Για τον \code{MarketplaceNFT}, του οποίου
το \en{Display} ορίζεται με την κλασική μορφή, η απόκριση ήταν κενή. Η εφεδρική ανάγνωση
από τα πεδία του αντικειμένου, που είχε προβλεφθεί εξαρχής, κάλυψε την περίπτωση χωρίς
αλλαγή στο συμβόλαιο.

\paragraph{Μορφοποίηση του κουμπιού σύνδεσης.} Το κουμπί σύνδεσης και το παράθυρο επιλογής
πορτοφολιού του \en{dApp Kit} είναι \en{web components} με \en{shadow DOM}, και δεν
μορφοποιούνται με συνήθεις κανόνες \en{CSS}. Προσαρμόζονται μόνο μέσω μεταβλητών που
κληρονομούν από τη σελίδα. Επειδή το παράθυρο αποδίδεται μέσα στο κουμπί, μια μεταβλητή
ακτίνας καμπυλότητας που έκανε το κουμπί στρογγυλό έκανε στρογγυλό και το παράθυρο.
Διατηρήθηκε μία κοινή τιμή και για τα δύο.

\paragraph{Εμφωλευμένα κουμπιά.} Στην πρώτη υλοποίηση, ολόκληρη η γραμμή μιας προσφοράς
ήταν ένα στοιχείο \code{button} που την ανέπτυσσε, και μέσα της βρίσκονταν τα κουμπιά
αποδοχής και απόρριψης. Η εμφώλευση αυτή δεν είναι έγκυρη \en{HTML} και, όπως έδειξε ο
αυτοματοποιημένος έλεγχος, καθιστούσε διφορούμενο το ποιο στοιχείο δέχεται το πάτημα. Στον
επανασχεδιασμό η κεφαλίδα της γραμμής έγινε στοιχείο με ρόλο κουμπιού και υποστήριξη
πληκτρολογίου, και τα κουμπιά ενεργειών είναι αδελφά στοιχεία.

\paragraph{Εξερευνητής και τοπικός κόμβος.} Ο σύνδεσμος της ειδοποίησης προς τον
\en{Polymedia Explorer} εμφάνιζε αρχικά «\en{Data could not be extracted}» για κάθε
συναλλαγή, παρότι ο κόμβος την επέστρεφε κανονικά. Ο εξερευνητής είναι σελίδα \en{HTTPS}
στο διαδίκτυο που, για το τοπικό δίκτυο, καλεί απευθείας τον κόμβο στη διεύθυνση
\code{127.0.0.1:9000} από τον φυλλομετρητή, και οι πρόσφατες εκδόσεις του \en{Chrome}
απαιτούν ρητή άδεια «πρόσβασης σε τοπικό δίκτυο» για τέτοιες κλήσεις. Το πρόβλημα δεν
έχει λύση από την πλευρά της εφαρμογής· η άδεια δίνεται μία φορά από τον χρήστη και
καταγράφεται στην τεκμηρίωση του αποθετηρίου.

\section{Επαλήθευση}
\label{sec:verification}

Η ορθή λειτουργία του συστήματος επαληθεύθηκε σε τρία επίπεδα, από το στενότερο προς το
ευρύτερο.

\begin{description}[leftmargin=1.2em,itemsep=3pt]
  \item[Έλεγχοι μονάδας του συμβολαίου.] Οι 21 έλεγχοι της Ενότητας~\ref{sec:tests}
    επαληθεύουν τη λογική των δύο \en{modules} σε προσομοιωμένο περιβάλλον.
  \item[Σενάρια συμβολαίου μέσω \en{SDK}.] Ένα σενάριο \en{TypeScript} δημιουργεί τρεις
    λογαριασμούς σε πραγματικό τοπικό δίκτυο, τους χρηματοδοτεί από το \en{faucet} και
    εκτελεί 20 σενάρια απευθείας στο συμβόλαιο, μετά από κάθε βήμα ελέγχοντας την
    κατάσταση της αλυσίδας (ιδιοκτησία, υπόλοιπα, ύπαρξη του \code{Bid}). Πέρα από τις
    κανονικές ροές και τους επτά κωδικούς σφαλμάτων, καλύπτει περιπτώσεις που δεν
    εκφράζονται εύκολα σε προσομοίωση: δύο προσφορές για το ίδιο \en{NFT} και τι συμβαίνει
    στη δεύτερη όταν γίνει αποδεκτή η πρώτη, δύο ταυτόχρονες αποδοχές της ίδιας προσφοράς
    από διαφορετικούς λογαριασμούς, προσφορά της οποίας ο στόχος είναι ο ίδιος δεσμευμένος
    σε άλλη προσφορά, αποδοχή της δικής του προσφοράς από τον προσφέροντα, ανταλλαγή
    \en{NFT} με νόμισμα μέσω των γενικών παραμέτρων, και λανθασμένα ορίσματα τύπου. Τα
    αποτελέσματα συζητούνται στην Ενότητα~\ref{sec:scenario-results}. Αυτό το επίπεδο
    αποκάλυψε το πρόβλημα της Ενότητας~\ref{sec:listevents}.
  \item[Αυτοματοποιημένη πλοήγηση της διεπαφής.] Με το \en{Playwright} \cite{playwright}
    ένας πραγματικός φυλλομετρητής ανοίγει την εφαρμογή, συνδέεται με το προσωρινό
    πορτοφόλι, και εκτελεί τις ροές πατώντας τα ίδια κουμπιά που θα πατούσε ένας χρήστης.
    Ένας δεύτερος και ένας τρίτος χρήστης, οδηγούμενοι από το \en{SDK}, παίζουν τον ρόλο
    των αντισυμβαλλομένων. Μετά από κάθε ενέργεια ελέγχεται το μήνυμα που εμφανίζεται και,
    ανεξάρτητα, η κατάσταση της αλυσίδας. Περιλαμβάνει 17 ροές, από τις σελίδες χωρίς
    σύνδεση και το \en{faucet} μέχρι την αποδοχή, την απόρριψη και την ακύρωση από
    οποιαδήποτε πλευρά, μια προσφορά που ακυρώθηκε ενώ ο χρήστης την έβλεπε, και τη
    διάταξη σε πλάτος 390 εικονοστοιχείων. Ένα συγγενές σενάριο παρήγαγε τα στιγμιότυπα
    του Κεφαλαίου~\ref{ch:epideixi}.
\end{description}

Η επιλογή του προσωρινού πορτοφολιού αξίζει μια εξήγηση. Τα πραγματικά πορτοφόλια είναι
επεκτάσεις φυλλομετρητή που ζητούν χειροκίνητη επιβεβαίωση, και ο αυτοματισμός τους είναι
εύθραυστος. Το προσωρινό πορτοφόλι υλοποιεί το ίδιο \en{Wallet Standard}, άρα η εφαρμογή
ακολουθεί ακριβώς την ίδια διαδρομή κώδικα, και επιτρέπει η πλοήγηση να φτάσει μέχρι την
αλυσίδα και πίσω χωρίς ανθρώπινη παρέμβαση. Ενεργοποιείται μόνο στο περιβάλλον ανάπτυξης.
